\paragraph{四重分块覆盖下的碰撞惯性提升、块核 \texorpdfstring{$3$}{3}-向量表示与局部半直积结构}
沿用本小节中 \(160\) 点分解 \(\bigsqcup_{j=1}^{40}B_j\)（每块 \(|B_j|=4\)）与块核 \(K\) 的记号，
并取碰撞边界的局部 braid 单值 \(\tau=\sigma_i^2\)（定理 \ref{thm:conclusion-s4-two-transposition-collision-three-type-boundary}）。
记块集合为 \(\Omega:=\{B_1,\dots,B_{40}\}\)，并以同一符号 \(\tau\) 表示其在 \(\Omega\) 上的诱导作用。

\begin{conclusion}[块纯提升与点态固定性]\label{con:conclusion-s4-collision-block-pure-lift-pointwise-fixedness}
在上述记号下：
\begin{enumerate}
  \item \(\tau\) 在 \(40\) 个块 \(\Omega\) 上的循环型为
  \[
  \tau\sim 1^{13}3^{9}.
  \]
  \item \(\tau\) 在 \(160\) 点纤维上的全部 \(3\)-循环均由上述 \(9\) 个块 \(3\)-循环轨道平行提升而来：每一个块 \(3\)-循环轨道在 \(160\) 点上提升为恰好四个互不相交的 \(3\)-循环。
  \item 被 \(\tau\) 固定的 \(13\) 个块皆为点态固定：若 \(\tau(B)=B\)，则 \(\tau|_{B}=\mathrm{id}\)。
\end{enumerate}
\end{conclusion}
\begin{proof}
由命题 \ref{prop:xi-terminal-zm-hurwitz-generator-cycle-types-on-160-and-40}，\(\sigma_i\) 在 \(\Omega\) 上循环型为 \(3^{9}1^{13}\)，且该表示中 \(|\sigma_i|=3\)；
因此 \(\tau=\sigma_i^2\) 在 \(\Omega\) 上仍为同一循环型，得第 (1) 点。
\par
由定理 \ref{thm:conclusion-s4-two-transposition-collision-three-type-boundary} 的证明，
\(\tau\) 在 \(160\) 点上满足
\[
\tau\sim 1^{52}3^{36}.
\]
取任一块 \(3\)-循环轨道 \((B_0,B_1,B_2)\)。
由于 \(\tau\) 将 \(B_k\) 双射地送到 \(B_{k+1}\)（指标模 \(3\)）且 \(\tau^3=1\)，对任意 \(x\in B_0\) 有
\[
x\ \mapsto\ \tau(x)\in B_1\ \mapsto\ \tau^2(x)\in B_2\ \mapsto\ \tau^3(x)=x,
\]
从而给出长度为 \(3\) 的点轨道。
若 \(x,y\in B_0\) 且 \(x\neq y\)，则其轨道不交（否则 \(y=\tau^k(x)\) 与 \(y\in B_0\) 矛盾），
故该块 \(3\)-循环轨道必提升为正好 \(|B_0|=4\) 个互不相交的点 \(3\)-循环。
因此 \(9\) 个块 \(3\)-循环轨道总计贡献 \(9\cdot 4=36\) 个点 \(3\)-循环，
与 \(\tau\) 在 \(160\) 点上的全部 \(3\)-循环数相同，得到第 (2) 点。
\par
由第 (2) 点知不存在任何“块内”\(3\)-循环。
于是 \(\tau\) 的 \(52\) 个不动点只能来自被 \(\tau\) 固定的 \(13\) 个块。
但每块大小为 \(4\)，故固定块贡献的不动点总数至多为 \(13\cdot 4=52\)，从而必处处取等：
每个固定块内四点皆不动，即 \(\tau|_{B}=\mathrm{id}\)。
得第 (3) 点。证毕。
\end{proof}

\begin{conclusion}[\texorpdfstring{$S_3$}{S3} 碰撞边界的不同指数与四倍律]\label{con:conclusion-s3-collision-different-exponent-fourfold-law}
设 \(f_{S_4}:\mathcal H_{S_4}\to \mathcal M\) 为定理 \ref{thm:conclusion-s4-two-transposition-collision-three-type-boundary} 的 \(S_4\)-Hurwitz 映射，
并令 \(f_{S_3}:\mathcal H_{S_3}\to \mathcal M\) 为相应 \(S_3\)-Hurwitz 映射；
记 \(\Delta\subset\mathcal M\) 为相邻分支点碰撞因子。
则：
\begin{enumerate}
  \item 正规化拉回 \(f_{S_3}^{-1}(\Delta)^\nu\) 具有 \(22\) 个几何不可约余维 \(1\) 分量：
  其中 \(13\) 个分量满足分歧指数 \(e=1\)，另有 \(9\) 个分量满足 \(e=3\)；
  因此
  \[
  \operatorname{diff}_{\Delta}(f_{S_3})=9\cdot(3-1)=18.
  \]
  \item 记 \(\pi:\mathcal H_{S_4}\to \mathcal H_{S_3}\) 为推论 \ref{cor:xi-terminal-zm-hurwitz-irreducible-and-genus2-jac3torsion-moduli} 所给的四重有限 \'etale 覆盖。
  则在 \(\Delta\) 的一般点处 \(\pi\) 不产生额外分歧，并且
  \(f_{S_4}^{-1}(\Delta)^\nu\) 的不可约分量在 \(\pi\) 下按 \(4{:}1\) 均匀分裂到
  \(f_{S_3}^{-1}(\Delta)^\nu\) 的分量之上，且分歧指数保持不变。
  因而
  \[
  \operatorname{diff}_{\Delta}(f_{S_4})=4\,\operatorname{diff}_{\Delta}(f_{S_3})=72,
  \]
  与推论 \ref{cor:conclusion-s4-boundary-different-exponent-72} 一致。
\end{enumerate}
\end{conclusion}
\begin{proof}
由命题 \ref{prop:xi-terminal-zm-hurwitz-generator-cycle-types-on-160-and-40}，
\(\sigma_i\) 在 \(40\) 点纤维（等价于块集合 \(\Omega\)）上的循环型为 \(3^{9}1^{13}\)，从而 \(\tau=\sigma_i^2\) 亦为 \(3^{9}1^{13}\)。
标准 Hurwitz 局部分歧理论表明：局部惯性在纤维上的每个长度 \(m\) 的循环对应归一化拉回边界上的一个不可约分量，且其分歧指数等于 \(e=m\)。
因此 \(f_{S_3}^{-1}(\Delta)^\nu\) 的分量数为 \(9+13=22\)，并有 \(9\) 个分量满足 \(e=3\)，得到
\(\operatorname{diff}_{\Delta}(f_{S_3})=\sum(e-1)=9\cdot 2=18\)。
\par
第 (2) 点中，\(\pi\) 在平滑部分为有限 \'etale；
并由结论 \ref{con:conclusion-s4-collision-block-pure-lift-pointwise-fixedness} 的块纯提升与点态固定性，
\(\tau\) 在 \(160\) 点纤维上的每个 \(3\)-循环都严格来自块 \(3\)-循环轨道的四条平行提升，而固定块内无任何内部 \(3\)-循环。
因此在 \(\Delta\) 的一般点处，\(\pi\) 在边界分量间仅作无分歧的 \(4{:}1\) 分裂，且保持分歧指数不变。
于是不同指数满足四倍律，
\(\operatorname{diff}_{\Delta}(f_{S_4})=4\,\operatorname{diff}_{\Delta}(f_{S_3})=72\)，并与推论 \ref{cor:conclusion-s4-boundary-different-exponent-72} 相符。证毕。
\end{proof}

\begin{conclusion}[\texorpdfstring{$\tau$}{tau} 在块核 \texorpdfstring{$3$}{3}-Sylow 上的 Jordan 型与幂零阶]\label{con:conclusion-s4-collision-tau-jordan-type-on-P3}
令 \(P_3\) 为定理 \ref{thm:conclusion-s4-block-kernel-sylow3-canonical-splitting} 的块核 Sylow--\(3\) 子群，
并置
\[
V:=P_3\otimes_{\ZZ}\FF_3.
\]
由同一定理的规范分解 \(P_3=\prod_{j=1}^{40}P_{3,j}\)（每个 \(P_{3,j}\cong C_3\)）可将 \(V\) 识别为 \(\FF_3^{40}\)。
令 \(\tau\) 由共轭作用诱导到 \(V\)，并记
\[
N:=\tau-\mathrm{Id}\in \End_{\FF_3}(V).
\]
则 \(N\) 的 Jordan 分解由 \(9\) 个大小为 \(3\) 的 Jordan 块与 \(13\) 个大小为 \(1\) 的 Jordan 块组成；
因而
\[
N^3=0,\qquad N^2\neq 0,\qquad
\operatorname{rank}(N)=18,\qquad
\operatorname{rank}(N^2)=9,\qquad
\dim_{\FF_3}\ker(N)=22.
\]
\end{conclusion}
\begin{proof}
由定理 \ref{thm:conclusion-s4-block-kernel-sylow3-canonical-splitting}，\(V\) 分解为 \(40\) 个一维坐标因子之直和，并由块坐标投影规范确定。
而 \(\tau\) 在块集合 \(\Omega\) 上的循环型为 \(3^{9}1^{13}\)（结论 \ref{con:conclusion-s4-collision-block-pure-lift-pointwise-fixedness}），
故 \(V\) 分解为 \(13\) 个不动坐标与 \(9\) 个长度为 \(3\) 的坐标轨道所张成的三维子空间之直和。
\par
在任一长度 \(3\) 的坐标轨道上，\(\tau\) 在某基下为一个（可能带单位标量的）三循环置换矩阵，
从而满足 \((\tau-\mathrm{Id})^3=0\) 且 \(\dim\ker(\tau-\mathrm{Id})=1\)；
因此该三维子空间对应唯一的长 \(3\) 幂零 Jordan 块。
在不动坐标上一维块中 \(N=0\)，对应 Jordan 块大小 \(1\)。
将 \(9\) 个三维 Jordan 块与 \(13\) 个平凡块逐块相加，即得 Jordan 型以及各秩与核维数。证毕。
\end{proof}

\begin{conclusion}[\texorpdfstring{$\langle\tau\rangle$}{<tau>} 在块核 \texorpdfstring{$3$}{3}-表示上的群上同调维数]\label{con:conclusion-s4-collision-c3-cohomology-dimension-formula}
令 \(C_3:=\langle\tau\rangle\) 为 \(\Delta\) 的局部惯性循环群（阶 \(3\)），并令 \(V\) 如结论
\ref{con:conclusion-s4-collision-tau-jordan-type-on-P3}。
则对一切 \(n\ge 0\) 有
\[
H^{2n+1}(C_3,V)\cong H^{2n}(C_3,V)\cong \FF_3^{13},
\]
特别地
\(
\dim_{\FF_3}H^1(C_3,V)=\dim_{\FF_3}H^2(C_3,V)=13
\)。
\end{conclusion}
\begin{proof}
由结论 \ref{con:conclusion-s4-collision-tau-jordan-type-on-P3}，\(V\) 在 \(\FF_3[C_3]\) 上分解为
\[
V\ \cong\ \FF_3^{13}\ \oplus\ J_3^{\oplus 9},
\]
其中 \(\FF_3\) 为平凡一维模，\(J_3\) 为三维不可分解 Jordan 模（\(\tau=\mathrm{Id}+N\)，\(N^3=0\)，且 \(\dim\ker N=1\)）。
\par
对特征 \(3\) 上的循环群 \(C_3\) 的标准周期分辨率给出
\[
H^{2n+1}(C_3,V)\cong \ker(\mathrm{Norm})/\mathrm{Im}(\tau-\mathrm{Id}),\qquad
H^{2n}(C_3,V)\cong \ker(\tau-\mathrm{Id})/\mathrm{Im}(\mathrm{Norm}),
\]
其中 \(\mathrm{Norm}=1+\tau+\tau^2\)。
在 \(\tau=\mathrm{Id}+N\) 且 \(\mathrm{char}(\FF_3)=3\) 时，
\(
\mathrm{Norm}=1+\tau+\tau^2=N^2
\)，从而
\[
H^{2n+1}(C_3,V)\cong \ker(N^2)/\mathrm{Im}(N),\qquad
H^{2n}(C_3,V)\cong \ker(N)/\mathrm{Im}(N^2).
\]
\par
在 \(J_3\) 上有
\(
\dim\ker(N)=1=\dim\mathrm{Im}(N^2)
\)
且
\(
\dim\ker(N^2)=2=\dim\mathrm{Im}(N)
\)，
故两类商空间均为零；
在平凡一维模上 \(N=0\)，两类商空间均为一维。
将 \(13\) 个平凡块求直和即得结论。证毕。
\end{proof}

\begin{conclusion}[\texorpdfstring{$\tau$}{tau} 的块轨道与局部 \texorpdfstring{$3$}{3}-群半直积结构]\label{con:conclusion-s4-collision-block-orbit-semidir-3group}
取 \(\tau\) 在块集合 \(\Omega\) 上的一条 \(3\)-循环轨道 \(O=\{B_0,B_1,B_2\}\)。
记
\[
P_3(O):=P_{3,B_0}\times P_{3,B_1}\times P_{3,B_2}\ \subset\ P_3,
\]
其中 \(P_{3,B}\) 表示与块 \(B\) 对应的一维因子（定理 \ref{thm:conclusion-s4-block-kernel-sylow3-canonical-splitting}）。
则：
\begin{enumerate}
  \item \(\tau\) 通过共轭在 \(P_3(O)\) 上诱导循环置换
  \[
  \tau P_{3,B_k}\tau^{-1}=P_{3,B_{k+1}}
  \qquad(\text{指标模 }3).
  \]
  \item 因而
  \[
  \langle\tau,\ P_3(O)\rangle\ \cong\ C_3^3\rtimes C_3,
  \]
  其中顶端 \(C_3\) 以 \(3\)-循环置换三个直因子；该群阶为 \(3^4=81\)。
  \item 令 \(X(O):=B_0\cup B_1\cup B_2\)（共 \(12\) 点），则 \(\langle\tau,P_3(O)\rangle\) 在 \(X(O)\) 上的置换表示恰有两个轨道：一个大小为 \(9\)，一个大小为 \(3\)。
\end{enumerate}
\end{conclusion}
\begin{proof}
第 (1) 点由 \(\tau\) 在块上的循环置换与 \(P_{3,B}\) 的块支撑性立即得到。
第 (2) 点为半直积结构的直接改写。
\par
对第 (3) 点，由定理 \ref{thm:conclusion-s4-block-kernel-sylow3-canonical-splitting}，
每个 \(P_{3,B_k}\) 在块 \(B_k\) 内为一个三循环并固定恰好一个点；
故 \(P_3(O)\) 在 \(X(O)\) 上固定三点（每块一固定点），并在其余九点上按块分解为三个 \(3\)-点轨道。
另一方面，\(\tau\) 循环置换三个块，因而将三固定点组成一个 \(3\)-点轨道；
并将九个非固定点在三块之间循环迁移，从而在这九点上生成单一轨道。
故轨道分解为 \(9+3\)。证毕。
\end{proof}

\begin{conclusion}[\texorpdfstring{$B$}{B}-型碰撞边界分支仅出现在 \texorpdfstring{$S_3$}{S3} 的 \texorpdfstring{$A$}{A}-型之上]\label{con:conclusion-s4-to-s3-btype-only-over-atype}
沿定理 \ref{thm:conclusion-s4-two-transposition-collision-three-type-boundary} 的碰撞边界记号，
令 \(p:=g_ig_{i+1}\) 表示相邻两换位的局部乘积。
在 \(S_4\) 中可出现三型 \(|p|\in\{1,2,3\}\)（对应 A/B/C 型），而在 \(S_3\) 中 \(|p|\in\{1,3\}\)。
对四重覆盖 \(\pi:\mathcal H_{S_4}\to\mathcal H_{S_3}\)（推论 \ref{cor:xi-terminal-zm-hurwitz-irreducible-and-genus2-jac3torsion-moduli}）有：
\begin{enumerate}
  \item 任一 C 型分支（\(|p|=3\)）在 \(\pi\) 下仍为 C 型；
  \item 任一 B 型分支（\(|p|=2\)）在 \(\pi\) 下必映到 A 型分支（\(|p|=1\)）。
\end{enumerate}
\end{conclusion}
\begin{proof}
记群商映射为
\(
\bar\pi:S_4\to S_4/V_4^{\mathrm{norm}}\cong S_3
\)。
其核为 \(V_4^{\mathrm{norm}}\)，非平凡元皆为阶 \(2\) 的双换位。
又 \(\bar\pi(g_i),\bar\pi(g_{i+1})\) 皆为 \(S_3\) 中的换位，故
\(\bar\pi(p)=\bar\pi(g_i)\bar\pi(g_{i+1})\) 只能为单位元或三循环，特别地不可能为阶 \(2\) 元。
因此若 \(|p|=2\)，则 \(\bar\pi(p)=1\)，像分支为 A 型，得第 (2) 点。
若 \(\bar\pi(p)\) 为三循环，则任一 lift 的阶必为 \(3\)（核无 \(3\)-部分），故 C 型在 \(\pi\) 下保持 C 型，得第 (1) 点。证毕。
\end{proof}

\begin{conclusion}[\texorpdfstring{$S_3$}{S3} 边界惯性在 \texorpdfstring{$S_4$}{S4} 四重提升中的“无内部扭曲”判别]\label{con:conclusion-s4-collision-no-internal-twist-uplift-criterion}
设 \(\tau=\sigma_i^2\) 为碰撞因子 \(\Delta\) 的局部惯性元。
则下述判别等价：
\[
\tau\ \text{在 \(160\) 点纤维上不存在任何块内 \(3\)-循环}
\quad\Longleftrightarrow\quad
\tau\ \text{的全部 \(3\)-循环均来自块 \(3\)-循环轨道的四条平行提升}.
\]
并且在等价情形下，对任意块 \(3\)-循环轨道 \(O\)，\(\tau\) 在
\(
P_3(O)\otimes_{\ZZ}\FF_3
\)
上的最小多项式必为 \((X-1)^3\)。
\end{conclusion}
\begin{proof}
由于 \(\tau\) 保持块系统且 \(|\tau|=3\)，任一点的块轨道长度只能为 \(1\) 或 \(3\)。
若存在块内 \(3\)-循环，则必存在某一固定块 \(B\) 使 \(\tau|_B\) 含 \(3\)-循环，
从而在 \(160\) 点循环分解中出现“固定块贡献的 \(3\)-循环”；
反之若不存在块内 \(3\)-循环，则所有 \(3\)-循环必来自块 \(3\)-循环轨道的提升。
\par
在后一情形下，由块间双射性与 \(\tau^3=1\)，每个块 \(3\)-循环轨道的提升必为四条互不相交的平行 \(3\)-循环；
从而右侧陈述成立。等价性得证。
\par
最小多项式结论由结论 \ref{con:conclusion-s4-collision-tau-jordan-type-on-P3} 的三维 Jordan 块描述立即推出：
在任一长度 \(3\) 的块轨道上，\(\tau\) 的不动子空间维数为 \(1\)，且 \((\tau-\mathrm{Id})^3=0\) 但 \((\tau-\mathrm{Id})^2\neq 0\)，故最小多项式为 \((X-1)^3\)。证毕。
\end{proof}

\begin{remark}
上述结论仅依赖于文稿中既定的循环型证书（命题 \ref{prop:xi-terminal-zm-hurwitz-generator-cycle-types-on-160-and-40} 与定理 \ref{thm:conclusion-s4-two-transposition-collision-three-type-boundary}）以及四重分块覆盖结构（结论 \ref{con:xi-terminal-zm-s4-to-s3-nielsen-40-blocks-of-4} 与定理 \ref{thm:conclusion-s4-block-kernel-sylow3-canonical-splitting}）。
因此其证明完全在离散层面封闭，可被视为后续几何与算术推导的结构性输入。
\end{remark}

